\documentclass[11pt,a4paper]{article}

\usepackage{fontspec}
\usepackage[spanish,es-tabla]{babel}
\usepackage[margin=2.2cm]{geometry}
\usepackage{booktabs}
\usepackage{longtable}
\usepackage{array}
\usepackage{xcolor}
\usepackage{enumitem}
\usepackage{tikz}
\usetikzlibrary{arrows.meta,positioning,babel}
\usepackage{graphicx}
\usepackage{hyperref}

\setlist{nosep,leftmargin=1.4em}

\newcommand{\costo}[1]{\par\smallskip\noindent\textbf{Costo estimado:} #1\par\bigskip}
\newcommand{\comp}[2]{\medskip\noindent{\large\textbf{#1}}\par\smallskip\noindent #2}

\title{\textbf{Sistema experimental completo: funcionamiento,\\
componentes, montaje y presupuesto}\\[4pt]
\large Reactor controlado para bioconversión con \emph{Hermetia illucens}\\
\large Versión 3.0}
\author{J.\,J. Leal --- tesis doctoral\\Repositorio: \texttt{github.com/asoleal/tesis-BSF}}
\date{Agosto 2026}

\begin{document}
\maketitle

\begin{abstract}
\noindent Este documento explica \textbf{cómo funciona el sistema completo}
(ruta del agua, del aire y de las señales), describe cada componente en un
párrafo claro ---qué es, cómo se instala y para qué sirve--- y presenta el
\textbf{presupuesto estimado}, incluyendo los costos de \textbf{ingeniería y
desarrollo de \emph{software}} si se contrata un ingeniero. Moneda de
referencia: USD; equivalencia aproximada 4\,000--4\,200~COP/USD (verificar al
comprar). Precios de referencia de mercado (agosto 2026), variación típica
$\pm$30\,\%.
\end{abstract}

\tableofcontents
\newpage

% ============================================================
\section{Cómo funciona el sistema (visión de conjunto)}

\begin{figure}[h!]
\centering
\resizebox{\textwidth}{!}{% Esquemático del sistema completo (vista de corte lateral)
\begin{tikzpicture}[font=\small, >=Stealth, line cap=round]

% ---------- colores ----------
\colorlet{agua}{cyan!15}
\colorlet{cama}{orange!25}
\colorlet{tapa}{gray!45}
\colorlet{aire}{red!80!black}

% ---------- tina del baño ----------
\draw[very thick] (0,0) rectangle (13.2,8.2);
\fill[agua] (0.08,0.08) rectangle (13.12,6.3);
\draw[blue, dashed] (0.08,6.3) -- (13.12,6.3);
\node[blue, anchor=west] at (0.15,5.55) {\scriptsize agua 30\,$^\circ$C};
\node[anchor=south east, gray] at (13.0,0.15) {\scriptsize tina (baño)};

% ---------- pilares ----------
\fill[gray!70] (2.4,0.25) rectangle (3.0,1.6);
\fill[gray!70] (10.0,0.25) rectangle (10.6,1.6);

% ---------- calentador ----------
\draw[fill=red!20] (12.0,0.8) rectangle (12.8,2.1);
\node[align=center] at (12.4,1.45) {\scriptsize calen-\\\scriptsize tador};

% ---------- celdas de carga ----------
\draw[fill=gray!50] (1.0,-0.75) rectangle (1.8,-0.1);
\draw[fill=gray!50] (11.4,-0.75) rectangle (12.2,-0.1);
\node[gray, align=center] at (6.6,-1.1) {\scriptsize 4 celdas de carga 50\,kg (peso continuo)};

% ---------- panera ----------
\draw[very thick] (1.6,1.6) rectangle (11.6,7.3);

% ---------- plenum ----------
\fill[gray!20] (1.7,1.7) rectangle (11.5,2.55);
\node at (6.6,2.1) {\scriptsize PLENUM (cámara de distribución)};
\draw[very thick, dash dot] (1.7,2.55) -- (11.5,2.55);
\node[anchor=west, gray] at (8.2,2.85) {\scriptsize placa difusora};

% ---------- bandejas (3 pisos) ----------
\foreach \ya/\yb in {2.85/3.75, 4.25/5.15, 5.65/6.55}{
  \draw[fill=cama] (2.1,\ya) rectangle (11.1,\yb);
  \draw[densely dotted, thick, brown] (2.1,\yb) -- (11.1,\yb);}
\foreach \x/\y in {2.6/3.1,3.5/3.35,4.5/3.05,5.6/3.35,6.7/3.1,7.7/3.3,8.8/3.05,9.8/3.3,10.6/3.1,
                   2.7/4.5,3.8/4.7,4.9/4.45,6.1/4.7,7.2/4.5,8.3/4.7,9.5/4.45,10.5/4.6,
                   2.6/5.9,3.6/6.1,4.7/5.85,5.8/6.1,6.9/5.9,8.0/6.15,9.2/5.85,10.3/6.1}{
  \fill[brown!80!black] (\x,\y) circle (0.055);}
\node[brown!70!black, anchor=east] at (1.55,3.3) {\scriptsize bandejas};
\node[brown, anchor=east] at (11.05,6.78) {\scriptsize geotextil};

% ---------- headspace ----------
\node[anchor=west] at (2.2,6.92) {\scriptsize headspace};

% ---------- tapa ----------
\draw[fill=tapa, very thick] (1.3,7.3) rectangle (11.9,7.72);
\draw[orange, very thick] (1.35,7.28) -- (11.85,7.28);
\node[anchor=west, gray] at (12.3,8.35) {\scriptsize tapa acrílica + empaque};
\draw[fill=black] (1.05,7.2) rectangle (1.35,7.95);
\draw[fill=black] (11.9,7.2) rectangle (12.2,7.95);
\node[anchor=west, gray] at (12.3,7.55) {\scriptsize abrazaderas};

% ---------- sensores en la tapa (etiquetas escalonadas) ----------
\foreach \x/\nom/\yl in {3.6/CO$_2$/8.45, 4.6/CH$_4$/8.95, 5.6/O$_2$/8.45, 6.6/T--HR/8.95}{
  \draw[fill=yellow!40] (\x-0.22,6.55) rectangle (\x+0.22,7.3);
  \draw[gray, thin] (\x,7.78) -- (\x,\yl-0.18);
  \node[anchor=south] at (\x,\yl-0.2) {\scriptsize \nom};
}
% septo
\draw[fill=blue!40] (8.6,7.3) rectangle (9.0,7.75);
\draw[gray, thin] (8.8,7.78) -- (8.8,8.28);
\node[anchor=south] at (8.8,8.28) {\scriptsize septo GC};
% válvula de salida
\draw[fill=magenta!30] (10.1,7.72) rectangle (10.7,8.15);

% ---------- sondas T de cama ----------
\foreach \x/\d in {9.5/4.35, 9.8/4.5, 10.1/4.85}{
  \draw[thick, purple] (\x,7.3) -- (\x,\d);
  \fill[purple] (\x,\d) circle (0.07);}
\node[purple, anchor=east] at (9.35,6.88) {\scriptsize T cama $\times$3};

% ---------- línea de aire de entrada (numerada) ----------
\foreach \i/\x in {1/-3.7, 2/-2.9, 3/-2.1, 4/-1.3, 5/-0.5}{
  \node[draw, circle, fill=red!10, inner sep=2.2pt] (c\i) at (\x,2.1) {\i};}
\draw[->, aire, thick] (c1) -- (c2);
\draw[->, aire, thick] (c2) -- (c3);
\draw[->, aire, thick] (c3) -- (c4);
\draw[->, aire, thick] (c4) -- (c5);
% tubería: atraviesa la pared de la tina y entra al plenum
\draw[->, aire, very thick] (c5.east) -- (1.78,2.1);

% ---------- flechas ascendentes por la cama ----------
\foreach \x in {3.3, 6.0, 8.4}{
  \draw[->, aire, very thick] (\x,2.2) -- (\x,6.45);}
\node[aire] at (3.62,3.05) {6.};
% salida por el headspace y la válvula
\draw[->, aire, very thick] (6.95,6.62) -- (10.35,6.62) -- (10.35,7.68);
\draw[->, aire, very thick] (10.4,8.18) -- (10.4,9.0);
\node[anchor=south, align=center] at (10.4,9.0) {\scriptsize 7. salida\\\scriptsize (válvula 3 vías)};

% ---------- electrónica ----------
\node[draw, fill=blue!10, minimum width=2.4cm, minimum height=1.1cm, align=center] (esp) at (15.2,5.2) {\scriptsize ESP32\\\scriptsize (control y\\\scriptsize adquisición)};
\node[draw, fill=blue!10, minimum width=2.4cm, minimum height=0.95cm, align=center] (rpi) at (15.2,3.4) {\scriptsize Raspberry Pi\\\scriptsize (dashboard / TSV)};
\draw[<->, thick, blue!60] (esp) -- (rpi);
\draw[dashed, gray, ->] (esp.west) -- (13.35,5.2);
\node[gray, align=center] at (14.6,6.25) {\scriptsize cables/Wi-Fi a\\\scriptsize sensores y válvula};
\draw[dashed, gray, ->] (rpi.south) |- (12.6,-0.5);

\end{tikzpicture}
}
\caption{Corte lateral del sistema. \textbf{Agua (celeste):} llena la tina y
rodea la panera (el tramo de tubería dentro del agua además precalienta el
aire). \textbf{Aire (flechas rojas, ruta numerada):} \textbf{1}~bomba,
\textbf{2}~filtro HEPA, \textbf{3}~humidificador de burbuja,
\textbf{4}~rotámetro, \textbf{5}~válvula solenoide, \textbf{6}~sube a través
de las bandejas, \textbf{7}~sale por la válvula de tres vías.
\textbf{Señales (líneas punteadas grises):} van y vienen del ESP32.}
\label{fig:esquema}
\end{figure}

\subsection{La idea en una frase}
Una \textbf{panera plástica} (el reactor) está \emph{sentada sobre pilares
dentro de una tina con agua caliente} que la abraza por fuera; adentro de la
panera hay \textbf{tres pisos de bandejas} con las larvas y su alimento; una
\textbf{bomba empuja aire por debajo} de las bandejas, el aire \emph{atraviesa
la cama de abajo hacia arriba} y sale por la tapa, donde unos \textbf{sensores}
miden los gases; una \textbf{tapa de acrílico con empaque} sella todo para que
ningún gas se escape sin ser medido; y un \textbf{ESP32} (un microcontrolador,
una pequeña computadora de USD~10) lee los sensores cada 30 segundos y abre o
cierra la válvula de aire según las reglas del gemelo digital.

\subsection{La ruta del agua (control de temperatura)}
El agua \textbf{no toca a las larvas}: queda \emph{por fuera} de la panera,
en la tina. Funciona como una ``chaqueta térmica'': como el agua guarda mucho
calor y cambia de temperatura lentamente, mantiene la panera a una temperatura
estable aunque la temperatura del cuarto suba o baje.

\begin{itemize}
  \item Un \textbf{calentador de inmersión con termostato} (o controlador
        STC-1000) mantiene el agua a 30~$^\circ$C de forma automática: el
        termostato apaga y prende el calentador sin intervención humana. Este
        es el \textbf{control de temperatura de base} (lento y estable).
  \item La panera se apoya en \textbf{pilares} para que el agua circule también
        por debajo, calentando parejo por todas las caras.
  \item Cuando las larvas comen mucho, \emph{ellas mismas generan calor} y la
        cama puede pasar de 33~$^\circ$C aunque el agua esté bien. Ahí actúa el
        \textbf{segundo nivel de control}: el ESP32 ordena más aireación y el
        aire fresco \emph{saca el calor} de la cama. Es decir: \textbf{el baño
        fija la base; la aireación quita los picos}.
\end{itemize}

\subsection{La ruta del aire (el corazón del experimento)}
El aire recorre un camino de una sola dirección (figura~\ref{fig:esquema}):

\begin{enumerate}
  \item \textbf{Bomba}: toma aire del cuarto.
  \item \textbf{Filtro HEPA}: quita polvo y esporas de hongos (para no
        contaminar el lote).
  \item \textbf{Humidificador de burbuja}: el aire atraviesa una botella con
        agua y sale húmedo. \emph{Este es el control de humedad del aire de
        entrada}: si el aire entrara seco, secaría la cama; entrando húmedo,
        la humedad del sustrato se mantiene en 65--70\,\%.
  \item \textbf{Rotámetro}: un tubo graduado con un flotador que muestra
        cuánto aire está pasando (lectura directa, sin electrónica).
  \item \textbf{Válvula solenoide}: una ``llave eléctrica'' que el ESP32 abre o
        cierra. Es la que ejecuta el control automático.
  \item El aire entra al \textbf{plenum} (cámara vacía bajo una placa
        perforada en el fondo de la panera) y \textbf{sube parejo a través de
        las tres bandejas}, llevándose el CO$_2$ y el calor que producen las
        larvas.
  \item Llega al \textbf{headspace} (el aire sobre la última bandeja), donde
        están los sensores, y \textbf{sale por la válvula de tres vías}:
        normalmente abierta (renovación continua) o cerrada 15 minutos al día
        (ventana de medición de flujos, protocolo P-05) --- al cerrarla, el
        gas se acumula y la pendiente de subida de CO$_2$ es la medida
        científica central de la tesis.
\end{enumerate}

\subsection{La ruta de las señales (quién manda sobre qué)}
Todos los sensores reportan al \textbf{ESP32} por cable (líneas punteadas de la
figura): CO$_2$, CH$_4$ y O$_2$ en la tapa, temperatura del aire y humedad
(SHT45), tres sondas de temperatura enterradas en la cama, humedad del
sustrato y el peso total (4 celdas de carga bajo la tina). El ESP32 decide la
aireación con una regla simple:

\begin{itemize}
  \item Si el CO$_2$ sube de 0,40\,\% $\rightarrow$ abre más la aireación.
  \item Si la cama pasa de 33~$^\circ$C $\rightarrow$ aireación al máximo
        (aunque el CO$_2$ esté bien).
  \item Si todo está en rango $\rightarrow$ aireación mínima (ahorra energía y
        no reseca la cama).
\end{itemize}

\noindent Los datos viajan por Wi-Fi a la \textbf{Raspberry Pi} (o al PC del
laboratorio), donde se guardan y se muestran en el \emph{dashboard} web.

\subsection{Resumen del control de temperatura y humedad}
\begin{center}
\begin{tabular}{@{}p{2.6cm}p{4.6cm}p{5.6cm}@{}}
\toprule
\textbf{Variable} & \textbf{Control de base} & \textbf{Control fino / de emergencia} \\
\midrule
Temperatura & Baño de agua a 30~$^\circ$C (termostato automático) &
Aireación forzada si la cama $>$ 33~$^\circ$C (ESP32 + solenoide) \\
Humedad del aire & Humidificador de burbuja en la entrada &
Bajar aireación mínima / rociar agua en la alimentación diaria (manual, registrado) \\
Humedad del sustrato & Preparación inicial a 65--70\,\% &
Sonda capacitiva avisa; ajuste con el agua de alimentación \\
\bottomrule
\end{tabular}
\end{center}

\newpage
% ============================================================
\section{Instalación del \emph{inlet} (entrada de aire), paso a paso}

\begin{enumerate}
  \item \textbf{Marcar}: un punto en la pared de la panera a 2~cm del fondo
        (por debajo de la placa difusora, es decir, \emph{dentro} del plenum).
        El aire debe entrar al plenum, no a la cama directamente.
  \item \textbf{Perforar}: orificio de 12~mm en la panera y otro alineado en
        la pared de la tina (a la misma altura, por encima del nivel del agua
        \emph{no es necesario}: la manguera puede entrar por arriba y bajar
        por fuera de la panera; así el agua nunca toca la unión).
        \textbf{Recomendado}: entrar \emph{por encima} --- la manguera pasa
        por el borde superior entre panera y tapa, bajando por dentro hasta el
        plenum; así no hay ninguna perforación bajo el nivel del agua.
  \item \textbf{Sellar en la panera}: racor pasamuros (\emph{bulkhead}) de
        6~mm con arandela de goma, o prensaestopa PG7 con silicona. Prueba de
        fugas P-01 (jabón en las uniones, cámara presurizada).
  \item \textbf{Conectar afuera, en este orden}: bomba $\rightarrow$ filtro
        HEPA $\rightarrow$ humidificador $\rightarrow$ rotámetro
        $\rightarrow$ solenoide $\rightarrow$ panera. Manguera de silicona de
        6~mm en todas las uniones, con abrazaderas plásticas.
  \item \textbf{Cablear}: la solenoide al relé del ESP32 (canal 1). El
        calentador del baño conserva su propio termostato (independiente).
  \item \textbf{Verificar}: con la tapa cerrada, abrir la bomba y confirmar
        burbujeo parejo bajo las tres bandejas (con la cama vacía y un poco de
        agua jabonosa sobre la placa difusora se ven las zonas activas).
\end{enumerate}

\newpage
% ============================================================
\section{Componentes (qué es, cómo va, para qué)}

\subsection{Reactor}

\comp{Panera plástica 34\,L (ya comprada)}
{Es el cuerpo del reactor: un recipiente rectangular de 38$\times$28~cm de
base y 26~cm de alto. No lleva modificación estructural: solo se apoya sobre
los pilares dentro de la tina con agua y recibe las perforaciones del inlet y
de los prensaestopas para los cables. Adentro van los tres pisos de bandejas
(${\approx}$4\,500 larvas) y queda el headspace donde se miden los gases.}
\costo{USD 0 (ya adquirida; referencia USD 8--15).}

\comp{Bandejas metálicas perforadas + geotextil (6 uds.)}
{Láminas de aluminio o acero inoxidable de ${\approx}$34$\times$24~cm con
perforaciones de 3~mm, forradas por encima con tela geotextil. Se apilan en
tres pisos dentro de la panera (dos bandejas por piso), separadas por topes de
${\approx}$10~cm. Los huecos dejan pasar el aire que viene del plenum; el
geotextil tapa los huecos para que el alimento y las larvas no se caigan. Se
fabrican en un taller local (corte y perforado de lámina).}
\costo{USD 10--15 por bandeja $\Rightarrow$ \textbf{USD 60--90}; geotextil
1~m$^2$: USD 5--10.}

\comp{Tapa rígida de acrílico + empaque + abrazaderas}
{Plancha de acrílico transparente de 45$\times$35~cm y 3--4~mm que cubre la
boca de la panera. Lleva pegado por debajo un cordón de empaque (goma EPDM o
silicona vulcanizada \emph{in situ}) y se presiona con 4 abrazaderas de
palanca (\emph{toggle}) atornilladas al borde de la panera. Sobre la tapa van
montados los sensores de gas, el septo y la salida. Sella herméticamente
---sin fugas, todo el gas medido proviene de la bioconversión--- y permite
abrir en segundos para alimentar o muestrear. Reemplaza la tapa rota.}
\costo{Acrílico USD 15--25; empaque USD 8--12; abrazaderas $\times$4 USD
10--16 $\Rightarrow$ \textbf{USD 33--53}.}

\comp{Placa difusora (plenum)}
{Una lámina perforada más (la misma de las bandejas) que se fija a 3~cm del
fondo de la panera, creando una cámara vacía debajo. El aire de la bomba
entra a esa cámara y sube uniformemente por todos los orificios: sin ella el
aire entraría como un chorro puntual y solo ventilaría una esquina.}
\costo{USD 10--15 (una lámina adicional del mismo pedido).}

\comp{Racores PG7, septo de GC y válvula de tres vías}
{Los PG7 son prensaestopas plásticos que sellan los cables que atraviesan la
pared (12~mm). El septo es un tapón de silicona autosellante pegado con epoxi
sobre un orificio de la tapa: permite meter la aguja de una jeringa y sacar
muestras de gas para el cromatógrafo sin abrir el reactor. La válvula de tres
vías en la salida dirige el aire a ``abierto'' (renovación continua) o
``cerrado'' (ventana de medición de 15 minutos).}
\costo{USD 15--25 el juego.}

\subsection{Aireación}

\comp{Bomba de aire 10~L/min}
{Bomba de diafragma tipo acuario, instalada fuera del baño. Empuja el aire
que atraviesa toda la línea hasta el plenum. Se dimensionó desde los datos
reales: con 6 bandejas (${\approx}$4\,500 larvas) se necesitan hasta
${\approx}$3,8~L/min; la bomba de 10~L/min da margen de sobra.}
\costo{USD 25--35.}

\comp{Filtro HEPA + humidificador de burbuja}
{El filtro (disco o cartucho HEPA) retiene polvo y esporas del aire que entra.
El humidificador es una botella con agua por la que el aire burbujea: sale
casi saturado de humedad y así el flujo continuo no reseca la cama. Van en
serie justo después de la bomba.}
\costo{USD 15--30.}

\comp{Rotámetro 0,2--4~L/min}
{Tubo transparente graduado con un flotador: muestra de un vistazo cuánto
aire está pasando. Va en serie antes de la entrada a la panera. Sirve para
fijar el caudal, registrarlo en la bitácora y detectar obstrucciones del
difusor.}
\costo{USD 20--40.}

\comp{Válvula solenoide 12~V + mangueras (10~m)}
{``Llave eléctrica'' del aire: el ESP32 la abre/cierra mediante un relé para
ejecutar el control automático de aireación. Las mangueras de silicona de
6~mm unen todos los elementos de la línea.}
\costo{USD 10--15 + USD 8--12.}

\subsection{Control térmico}

\comp{Tina + calentador de inmersión con termostato}
{La tina es cualquier caja plástica mayor que la panera; el calentador (tipo
acuario, 100--200~W, con termostato propio o comandado por un STC-1000)
mantiene el agua a 30~$^\circ$C automáticamente. La panera flota/apoya sobre
pilares para que el agua la abrace por debajo y los lados. Es el control de
temperatura de base del sistema.}
\costo{USD 25--40; tina USD 0--15 (reutilizable); pilares USD 5.}

\comp{Cobertura oscura}
{Tela opaca o caja que cubre el conjunto: las larvas requieren oscuridad y
además estabiliza un poco más la temperatura ambiente alrededor del equipo.}
\costo{USD 10--20.}

\subsection{Sensores}

\comp{CO$_2$ NDIR Winsen MH-410D 0--5\,\%vol (compra nueva)}
{Sensor infrarrojo de dióxido de carbono, la medida central de la tesis. Se
atornilla a la tapa con el cabezal dentro del headspace y el cable al ESP32
por un PG7. La versión actual del laboratorio (0--5\,000~ppm) saturaba durante
los ensayos; esta llega a 50\,000~ppm y elimina el problema de enmascaramiento
que afectó a los lotes 9--17.}
\costo{USD 60--120.}

\comp{CH$_4$ NDIR Winsen MH-441D (ya comprado)}
{Sensor infrarrojo de metano, montaje idéntico al de CO$_2$. Mide el metano
del lote y se evalúa como señal temprana de cosecha. Pendiente la validación
contra GC con gas real del batch (sensibilidad cruzada a VOCs).}
\costo{USD 0 (adquirido; referencia USD 50--120).}

\comp{O$_2$ óptico SST LuminOx LOX-02 (compra nueva)}
{Sensor de oxígeno por fluorescencia: 0--25\,\%, exactitud $<$2\,\% FS, vida
$>$5 años, no se consume como los electroquímicos. Se monta en la tapa junto a
los otros dos sensores de gas. Con O$_2$ y CO$_2$ simultáneos se calcula el
cociente respiratorio (RQ), que separa respiración larvaria de fermentación
microbiana: información clave para el gemelo y novedad frente a la literatura.}
\costo{USD 90--150 (el ítem más costoso).}

\comp{SHT45 (2 uds.) --- temperatura y humedad del aire}
{Módulo digital de alta precisión ($\pm$0,1~$^\circ$C, $\pm$1,5\,\%HR). Uno en
el headspace (condiciones del gas medido) y otro en el ambiente del
laboratorio. Permiten corregir las lecturas de ppm a condiciones reales.}
\costo{USD 12--15 c/u $\Rightarrow$ USD 24--30.}

\comp{DS18B20 sumergibles (3 uds.) --- temperatura de la cama}
{Sondas digitales encapsuladas en acero, enterradas a 1, 2,5 y 4~cm en la
bandeja central. Miden el perfil vertical de temperatura: detectan el
sobrecalentamiento interno que motivó este rediseño y validan el modelo
térmico del gemelo.}
\costo{USD 3--5 c/u $\Rightarrow$ USD 10--15.}

\comp{Humedad de sustrato capacitiva (2 uds.)}
{Sondas que van enterradas en la cama de dos bandejas y reportan la humedad
del sustrato (meta: 65--70\,\%). La versión capacitiva no se corroe como las
resistivas.}
\costo{USD 3--5 c/u $\Rightarrow$ USD 6--10.}

\comp{Celdas de carga 50~kg (4 uds.) + HX711}
{Cuatro galgas de báscula bajo las esquinas de la tina y un convertidor de 24
bits. Pesan el lote completo en continuo, sin abrir el reactor: la curva de
peso es una de las señales de cosecha y calibra el crecimiento larvario del
gemelo.}
\costo{USD 12--20 el kit.}

\subsection{Adquisición y cómputo}

\comp{ESP32 DevKit}
{Microcontrolador con Wi-Fi (una pequeña computadora de propósito único). Lee
todos los sensores cada 30~s, guarda los datos (TSV), ejecuta la regla de
control de aireación y reporta a la Raspberry Pi. Va en una caja eléctrica
fuera del baño.}
\costo{USD 8--15.}

\comp{Raspberry Pi (o PC reutilizado)}
{Computadora del tamaño de una tarjeta que recibe los datos por Wi-Fi, los
almacena con respaldo y sirve el \emph{dashboard} web del laboratorio.
Alternativa de costo cero: usar el PC existente.}
\costo{USD 60--80 (o USD 0).}

\comp{Fuente 5~V/5~A, relés, cableado, protoboard}
{Alimentación estable y módulo de relés para comandar la solenoide (y
opcionalmente el calentador).}
\costo{USD 30--40.}

\newpage
% ============================================================
\section{Desarrollos de ingeniería y \emph{software}}

El \emph{software} es libre (sin licencias), pero el desarrollo toma tiempo de
ingeniería. Dos escenarios: \textbf{(a) desarrollo propio} (usted, como parte
de la tesis) o \textbf{(b) contratar un ingeniero} (electrónica/
programación; tarifa de referencia Colombia: USD 8--20/h según experiencia;
estudiante en práctica o semillero puede bajar al mínimo).

\begin{longtable}{@{}>{\raggedright\arraybackslash}p{5.2cm}>{\raggedright\arraybackslash}p{6.2cm}cc@{}}
\toprule
\textbf{Tarea} & \textbf{Qué incluye} & \textbf{Horas} & \textbf{Costo (USD)} \\
\midrule
\endhead
Montaje del reactor y cableado & perforaciones, sellos, línea de aire,
caja eléctrica, conexión de sensores y relés & 40--60 & 320--1\,200 \\
Firmware ESP32 & lectura de buses (UART/I$^2$C/1-Wire/ADC/HX711), TSV cada
30\,s, regla de control ACH, Wi-Fi & 40--80 & 320--1\,600 \\
\emph{Dashboard} web & Panel/Dash en la Raspberry Pi: gráficas en vivo,
alarmas, vista gemelo vs.\ medido, señal de cosecha & 60--100 & 480--2\,000 \\
Puesta en marcha y calibración conjunta & prueba de fugas, calibración de
celdas, verificación cruzada de sensores, primer batch de prueba & 20--40 & 160--800 \\
\midrule
\textbf{Total contratado} & & \textbf{160--280} & \textbf{1\,280--5\,600} \\
\bottomrule
\end{longtable}

\noindent\textbf{Recomendación:} el gemelo digital y la calibración PINN los
hace usted (son el núcleo de la tesis y ya van avanzados); el montaje y el
firmware pueden contratarse parcialmente o hacerse con apoyo de un semillero.
Escenario mixto estimado: \textbf{USD 500--1\,500}.

% ============================================================
\section{Análisis de laboratorio}

\comp{Caracterización de los alimentos (análisis proximal + C/N + pH)}
{Se envían muestras de ${\approx}$300~g de cada alimento (sustrato de pollo y
el segundo tipo, en los días 0, 3 y 6 de fermentación) a un laboratorio de
análisis de alimentos: humedad, cenizas, proteína, fibra, grasa, relación C/N
y pH. El gemelo necesita la composición real del sustrato para predecir gases
según el alimento, y es requisito de reporte en revistas Q1.}
\costo{Proximal USD 50--100 por muestra + C/N y pH USD 20--40; con 4--6
muestras: \textbf{USD 280--600} (menos si es laboratorio de la facultad).}

\comp{Validación de gases por cromatografía (GC) + gas patrón}
{Muestras tomadas con jeringa por el septo (minuto 12 de la ventana cerrada)
inyectadas al GC, más prueba de interferencia con etanol, y un gas patrón
certificado para chequeo periódico. El CO$_2$ ya se validó ($\pm$20~ppm); el
CH$_4$ requiere este paso para confirmar que lo medido es metano y no VOCs.}
\costo{USD 10--30 por inyección ($\approx$20 muestras: USD 200--400) + gas
patrón USD 50--150 $\Rightarrow$ \textbf{USD 250--550} (USD 0--150 con
convenio universitario).}

\newpage
% ============================================================
\section{Resumen de presupuesto}

\begin{longtable}{@{}>{\raggedright\arraybackslash}p{8.6cm}rr@{}}
\toprule
\textbf{Bloque} & \textbf{Mín. (USD)} & \textbf{Máx. (USD)} \\
\midrule
\endhead
Reactor (bandejas, tapa, difusor, racores) & 118 & 183 \\
Aireación (bomba, filtro, humidif., rotámetro, solenoide) & 78 & 132 \\
Control térmico (tina, calentador, cobertura) & 40 & 75 \\
Sensores (compras nuevas) & 212 & 350 \\
Adquisición y cómputo (ESP32, RPi, fuente, relés) & 98 & 135 \\
Análisis de laboratorio (alimentos + GC) & 530 & 1\,150 \\
\midrule
\textbf{Subtotal materiales y laboratorio} & \textbf{1\,076} & \textbf{2\,025} \\
Ingeniería y desarrollo \textbf{contratado} & 1\,280 & 5\,600 \\
\quad (escenario mixto recomendado) & (500) & (1\,500) \\
\midrule
\textbf{Total si todo se contrata} & \textbf{2\,356} & \textbf{7\,625} \\
\textbf{Total escenario mixto (recomendado)} & \textbf{1\,576} & \textbf{3\,525} \\
\textbf{Total si usted desarrolla todo} & \textbf{1\,076} & \textbf{2\,025} \\
\bottomrule
\end{longtable}

\noindent\textbf{Cifra de planeación recomendada: ${\approx}$USD 2\,000--2\,500}
(${\approx}$COP 8--10 millones a 4\,100~COP/USD), que cubre materiales,
laboratorio y apoyo parcial de ingeniería, con 15\,\% de imprevistos.

\medskip
\noindent\textbf{Compras que bloquean el avance (hacer primero):}
\begin{enumerate}
  \item MH-410D 0--5\,\%vol --- sin él, el CO$_2$ satura otra vez.
  \item Bandejas metálicas + geotextil y tapa de acrílico --- sin ellas no hay
        reactor hermético ni aireación de fondo.
  \item Bomba 10~L/min + rotámetro + solenoide --- control de aireación.
  \item LOX-02 --- puede comprarse en segunda fase, antes de la campaña
        definitiva.
\end{enumerate}

\noindent\textbf{Ahorros posibles:} reutilizar PC ($-$USD~80), GC y gas patrón
por convenio ($-$USD~550), proximal en la facultad ($-$USD~300), desarrollo con
semillero/práctica ($-$USD~800--4\,000).

\section*{Notas}
\begin{itemize}
  \item Precios de referencia consultados en agosto 2026: MH-410D
        USD~55--120 según proveedor; LOX-02 USD~60--145; SHT45 USD~12,5--15;
        tarifas de ingeniería Colombia USD~8--20/h (junior--medio).
  \item Los costos de laboratorio son cotizaciones típicas; solicitar
        cotización formal antes de la campaña.
  \item Ítems ya adquiridos (paneras, MH-441D, MH-410D 5\,000~ppm) figuran con
        costo USD~0 y quedan como respaldo o controles.
\end{itemize}

\end{document}
